\scanpage{155}
% archive foliation: 154
ramener au cas où $G$ opère transitivement sur l'ens des comp. connexes.
Soit $X_0$ \inserted{$(0\in I)$} une d'elles, $G_0$ son stabilisateur, alors il est évident que $X$ s'identifie à $\unsure{X|X_0}\times_{G_0}G$, et $X/G$ à $X_0/G_0$.
\deleted{Enfin, \unsure{donc}} on est ramené à $X_0$ connexe.
\deleted{Soit $H$ \unsure{$x\in X_0$}, soit $G_i$ son groupe d'inertie, il opère trivialt sur $X$ \unsure{puisque est connexe}.}
\deleted{\unsure{l'ens}}
On peut supposer que $G_0$ opère \uline{\unsure{fidelt}}.
Mais \uline{alors il opère sans inertie} donc $X\longrightarrow X/G$ est \uline{étale} (prop ---).

Reste à voir que $X/G$ est étale sur $Y$.

\margin{$*\quad\unsure{\prime}$}
\diagram{\begin{tikzcd}
X \arrow[dr] & \\
& X'=X/G \arrow[d] \\
& Y\ni y
\end{tikzcd}
\qquad
\begin{tikzcd}
k(x') \arrow[d,no head] \\
k(x) \arrow[d,no head] \\
k(y)
\end{tikzcd}}

Faisons une extension fid plate de la base $Y$ convenable, on est ramené au cas où l'ext résiduelle $k(x)/k(y)$ est triviale.
\unsure{Comme}, on voit sur le diagramme $\widehat{\mathcal O}_y\simeq\widehat{\mathcal O}_{x'}$, cqfd.

\diagram{\begin{tikzcd}[column sep=large]
\widehat{\mathcal O}_y \arrow[r,no head] \arrow[rr,bend right=25,"\sim"'] &
\widehat{\mathcal O}_{x'} \arrow[r,"\sim"] &
\widehat{\mathcal O}_x
\end{tikzcd}}

\uline{Corollaire}
\deleted{Soit} $X$ étale sur $Y$ en $x$, $q:X\to X'$ un $Y$-morphisme étale en $x$; alors $X'$ est étale sur $Y$ en $q(x)$.

\figure{\begin{tikzpicture}[x=1cm,y=1cm]
\draw (2.8,5.2) .. controls (1.9,4.8) and (2.6,3.8) .. (2.8,3) .. controls (3,2.1) and (3,1.2) .. (2,1) -- (1.4,1.2) .. controls (.1,.9) and (-.2,.2) .. (.4,0) .. controls (1.4,-.3) and (2.2,.1) .. (2.1,.9);
\end{tikzpicture}}
